% Fig 3: cut-edge mechanism. Ensemble cut-edge distribution per state (grey
% histogram) with the enacted plan marked (orange line). The field-standard
% ensemble-validity / compactness figure, doubling as the "why" for Fig 2.
\begin{figure}[t]
\centering
\begin{tikzpicture}
\begin{groupplot}[
  urtc, height=2.7cm, width=\columnwidth,
  group style={group size=1 by 3, vertical sep=7mm},
  ytick=\empty, ylabel={},
  every axis title/.style={at={(0.02,0.94)}, anchor=north west, font=\scriptsize},
]
\nextgroupplot[title={NC: enacted in the bulk (57th pctile)}, xmin=600, xmax=945,
  title style={at={(0.98,0.94)}, anchor=north east, font=\scriptsize}]
\addplot+[hist={bins=22}, fill=cMUTE!55, draw=cGREY!70, mark=none, forget plot]
  table[y=cut_edges, col sep=comma]{ensemble_NC.csv};
\draw[cORG, very thick] (axis cs:749,0) -- (axis cs:749,200);

\nextgroupplot[title={PA: enacted past the tail ($+5.8\sigma$)}, xmin=1500, xmax=2450]
\addplot+[hist={bins=22}, fill=cMUTE!55, draw=cGREY!70, mark=none, forget plot]
  table[y=cut_edges, col sep=comma]{ensemble_PA.csv};
\draw[cORG, very thick] (axis cs:2367,0) -- (axis cs:2367,200);

\nextgroupplot[title={MD: enacted past the tail ($+13\sigma$)}, xmin=230, xmax=830, xlabel={cut edges (enacted plan, orange)}]
\addplot+[hist={bins=22}, fill=cMUTE!55, draw=cGREY!70, mark=none, forget plot]
  table[y=cut_edges, col sep=comma]{ensemble_MD.csv};
\draw[cORG, very thick] (axis cs:777,0) -- (axis cs:777,200);
\end{groupplot}
\end{tikzpicture}
\caption{The cut-edge mechanism. Grey: the distribution of cut edges (dual-graph
edges crossing a district boundary, the discrete compactness statistic a GIN is
sensitive to) over each state's neutral ensemble. Orange: the enacted plan. North
Carolina's enacted plan sits in the body of its distribution, so the geometry
detector correctly sees nothing anomalous; Pennsylvania's and Maryland's lie far
beyond their ensemble tails, which is why the topology-only detector flags them.}
\label{fig:cutedges}
\end{figure}
